\section{Life as an Engineering Property}

\subsection{An operational definition}

Biological definitions of life disagree at the margins, while computational
systems can reproduce selected properties without reproducing biology. We
therefore avoid a binary metaphysical definition and specify a functional
profile. A system is a \emph{machine proto-organism} to the degree that it
demonstrates all of the following:

\begin{enumerate}[leftmargin=*]
  \item \textbf{Boundary:} it distinguishes internal state, authorized
  components, and admitted knowledge from external resources and proposals.
  \item \textbf{Continuity:} its identity and direction survive replacement
  of individual processes, models, sessions, and machines.
  \item \textbf{Metabolism:} it acquires compute, storage, code, toolchains,
  and information, converting them into maintained executable capacity.
  \item \textbf{Self-maintenance:} it detects degradation and performs
  bounded repair, update, backup, and recovery work.
  \item \textbf{Adaptation:} experience changes future behavior through
  inspectable memory, code, policy, or organ revisions.
  \item \textbf{Action:} it can cause effects in its habitat under explicit,
  revocable authority.
  \item \textbf{Self-production:} it can participate in producing qualified
  successor versions of parts of its own organization.
\end{enumerate}

No one property is sufficient. A cron job has recurrence but little
continuity. A language model has general inference but no durable boundary. A
repository has heredity but no action. A self-updating service has
self-modification but may have no inspectable direction or identity.

The proposal is also graded. A laboratory system may meet these properties
weakly and depend on a human-supplied founding Initiative, purchased compute,
and externally maintained power. Biological organisms likewise depend on
environmental conditions and inherited drives. The engineering question is
not whether dependence disappears, but whether the system organizes that
dependence into sustained, bounded activity.

\subsection{Subject identity is not process identity}

Let a subject at admitted state cut \(S_t\) be:

\[
S_t = (I, D_t, F_t, A_t, H_t, O_t, B_t)
\]

where \(I\) is stable subject identity, \(D_t\) is continuing direction,
\(F_t\) is the admitted fact root, \(A_t\) is current authority, \(H_t\) is
eligible causal history, \(O_t\) is executable organization, and \(B_t\)
describes the habitat boundary. An agent invocation \(g_t\) reads a bounded
view of \(S_t\), proposes work, and may produce candidate evidence. It does not
become \(S_t\), and its output does not become \(S_{t+1}\) merely because it
was generated.

A successor transition requires an admission function:

\[
S_{t+1} =
\begin{cases}
\operatorname{admit}(S_t, C_t, E_t, W_t), & \text{if qualification succeeds},\\
S_t, & \text{otherwise},
\end{cases}
\]

where \(C_t\) is a candidate change, \(E_t\) is evidence, and \(W_t\) is the
warrant authorizing the effect. This makes semantic continuity compatible
with material replacement. Models can change, sessions can terminate, caches
can be rebuilt, and machines can be replaced without silently changing the
identity or declared direction of the subject.

\subsection{Needs and direction}

A primitive organism need not invent its founding need. A human can establish
an Initiative to develop a safer and more capable machine organism without
exceeding the admitted authority boundary. What matters is not the origin of
the sentence but its durable role: it must outlive individual actions,
constrain successor work, be revisable only through explicit authority, and
generate concrete maintenance and development pressure.

\inferencelabel\ Given a continuing Initiative, recurrent scheduling, access
to a sufficiently capable agent, and authorized engineering tools, the system
can exhibit strong directionality without continuous human attention.

\nonclaimlabel\ A supplied Initiative does not prove intrinsic desire,
conscious intention, affect, or independent moral agency. The term
\emph{need} in this paper denotes a persistent control constraint that
allocates work and resources.

\subsection{Semantic autopoiesis}

Varela, Maturana, and Uribe described autopoietic organization in terms of a
network that produces components which recursively realize the network and
its boundary \cite{varela1974autopoiesis}. The present proposal does not claim
that software thereby becomes biologically autopoietic. It extracts an
engineering analogue:

\begin{quote}
\emph{Semantic autopoiesis is the governed production of successor executable
organization by a system whose identity, direction, fact boundary, authority,
and admission rules persist across the production process.}
\end{quote}

The adjective \emph{semantic} is load-bearing. If a system can rewrite code but
cannot say which subject authorized the rewrite, which facts justified it,
which direction it serves, or whether the result was admitted, it has
mutation without accountable self-production. Conversely, a perfect semantic
record without executable regeneration has continuity without metabolism.
The proposed organism requires both.
